\section{Introduction}\label{sec:introduction}

Researchers who study long documents (party manifestos, legislative records, clinical narratives) routinely need a single judgment that depends on information spread across the entire text. An expert RILE score for a manifesto reflects how environmental, economic, and social positions combine; a clinical adverse-event count depends on interactions between symptoms reported pages apart. Producing these judgments at scale requires either expensive full-document annotation or some form of compression. C-TreePO provides the compression and proves when it is safe.

The idea is simple. Split a long document into contiguous spans, summarize each span, merge sibling summaries pairwise up a tree, and predict the task value from the root. If the summaries preserve the right information at each step, the root prediction is as good as reading the full document. We formalize ``the right information'' as three local laws (sufficiency, idempotence, and merge consistency) and show that when they hold, the global task signal is preserved. The proof object is a triangle. Any span in the tree can be written three ways: as raw text, as a summary produced by one call on the whole span, or as a merge of its two halves' summaries. Merge consistency demands that all three carry the same task value, and stacking these triangles from the leaves upward carries that agreement to the root. This is the recipe of a \emph{mergeable sketch} from the streaming-data-systems literature \citep{Agarwal2013MergeableTODS}, generalized from hand-designed merge rules to learned neural summarizers; Section~\ref{sec:mergeable-sketches} makes the analogy precise and Appendix~\ref{sec:hll-parity} shows the framework recovering HyperLogLog byte-for-byte.

A concrete example shows what can go wrong without the local laws. Consider a $128$-token document generated by a $4$-color Markov chain, where the task is to count how many times the hidden color changes. If we split the document into contiguous leaves and process each leaf independently, a color transition at a leaf boundary is invisible to either leaf. A naive sum of per-leaf counts systematically undercounts. The same failure occurs in semantic settings when a merge summary drops a cross-chapter tension (say, that a manifesto is environmentally interventionist but economically restrained) that determines the full-document judgment. The merge must carry enough information to detect what happens at the boundary.

The framework offers three things. First, preservation: if the local laws hold, the tree's root summary supports the same task judgment as the full document, including for DPO, GRPO, and related preference optimizers. Second, an operational audit: the local laws can be checked by sampling a small number of nodes and querying a local oracle, without re-reading the full document. Third, a finite-sample gap bound that decomposes into an empirical IPW estimate, a fixed calibration residual (the gap between the judge and the task oracle, independent of label count), and an estimation envelope that shrinks monotonically in the total number of labels collected, whether global (document-level) or local (node-level). The learning view is just as simple: fit the task oracle, and use local-law weights as a projection toward the theorem-eligible subspace. Equivalently, the approximation error is structured so that zero local-law error means exactly C1/C2/C3 validity. Section~\ref{sec:theorems} states this projection/structured-error equivalence in the main text, and Appendix~\ref{app:projection-interpretation} gives the formal crosswalk.

The preservation guarantee is about equality: when the local laws hold, the tree's prediction matches the full-document oracle. The practical value comes from a second observation: the same local signal the audit samples is a signal the model can also be supervised on. When the task's structure aligns with the local laws (additive oracles, boundary-sensitive interactions, rubric-coded spans), local labels pay for root accuracy at a lower total cost than global labels alone. Appendix~\ref{sec:markov-walkthrough} shows this directly on the Markov benchmark: roughly half the root labels can be replaced by local labels at no cost to root accuracy, and the tree matches or improves on the flat baseline across the full budget range.

The centerpiece application is political-text measurement on the
Manifesto Project corpus, at two levels of supervision. With
document-level labels only (the expert-survey means of
\citealp{BenoitEtAl2025}), a prompt-only C-Tree over a single
open-weight model matches a three-frontier-LLM ensemble (macro Pearson
$0.829$ vs.\ $0.817$) at roughly two orders of magnitude lower
compute, and its root score moves by only $0.027$ while the leaf size
varies $16\times$: the chunk schedule is not the estimand. The
deployed objects there are prompt programs (a summarizer/merge prompt
and a scorer prompt), so Section~\ref{sec:rile-fg-ladder} also asks
whether repeated prompt updates to the two roles preserve the same
divide-and-conquer behavior, rather than assuming that a single
full-context read or a one-time prompt choice settles the problem.
The same corpus then removes the prompt-program complication entirely:
the Manifesto Project's expert quasi-sentence codes aggregate
additively over spans, so every node of every tree has a gold oracle
value, and Section~\ref{sec:manifesto-qsentence} trains both a
prompt-program substrate and a neural-operator substrate against that
supervision on identical bundles, with the local laws entering as
training signals rather than audit probes.

\subsection{The Object We Preserve}

The object of interest is an \emph{oracle} judgment $\fstar(x) \in \Y$: the task value a domain expert would assign by the expensive, full-information procedure the researcher trusts. In a controlled setting this may be an exact count---the number of color transitions in a token sequence, or the number of adverse events in a clinical record. In political measurement it may be an expert RILE score. The summarizer $g$ discards text irrelevant to $\fstar$ while retaining the distinctions that $\fstar$ needs. The summary is a task-specific sufficient statistic for the oracle value.

This framing connects the algebraic literature on mergeable state with the
measurement literature on surrogate labels. Section~\ref{sec:mergeable-sketches}
develops the first link: a safe compression is a task-relative state whose
merge preserves the readout. Section~\ref{sec:social-measurement} develops
the second: in social-science applications, the readout is an expert
measurement target, and replacing full-document annotation by cheaper model
calls is useful only when the downstream target is preserved.

Manifesto measurement is the running reference; the Markov changepoint task is the controlled anchor. In the controlled task, documents are $128$-token sequences generated by a $4$-color Markov chain and the oracle $\fstar$ is the number of color transitions. Appendix~\ref{sec:markov-interlude} motivates why unordered symmetric merges without boundary state cannot handle that task, Appendix~\ref{sec:fno-primer} introduces the equation-(6) neural-operator realizability layer, its transformer architecture inclusion surface, and the FNO instance used in the controlled benchmark, and Appendix~\ref{sec:markov-walkthrough} shows the tree matching or improving on a flat baseline across the full range of supervision budgets. The same mechanism collapses to classical mergeable sketches at one extreme (Appendix~\ref{sec:hll-parity} recovers HyperLogLog byte-for-byte) and carries a full semantic load at the other (Section~\ref{sec:manifesto-deepdive}).

\subsection{Result Ladder}

The technical story is a nested ladder. The neural-operator layer supplies a realizable operator on the compact set of calls the tree actually makes. The oracle/projection objective then trades oracle fit against local-law projection; the projection/structured-error equivalence says this is the same assumption as requiring approximation error to vanish exactly on the local-law subspace. Those budgets feed the existing preservation and gap theorems.
\begin{enumerate}[leftmargin=1.5em]
    \item \textbf{Local laws $\rightarrow$ root preservation.} Under the \emph{Preservation Stack} (C1, C3, and context compatibility), local leaf and merge checks compose to root-level oracle preservation and schedule invariance. Adding C2 yields safe extra re-summarization rounds.
    \item \textbf{Preserved oracle $\rightarrow$ same training target.} Under the \emph{Optimization Stack}, which adds factorization and oracle-measurability assumptions, DPO, GRPO-PL, and GRPO-RL on summaries target the same population argmin set as training on full documents.
    \item \textbf{Approximate preservation $\rightarrow$ empirical certificate.} Under the \emph{Certificate Stack}, which adds transport, calibration, and sampling assumptions, sampled node audits yield a finite-sample bound on preference-objective drift.
\end{enumerate}

On the Markov changepoint benchmark, decomposing the document into a tree of subproblems matches or improves on a flat Fourier Neural Operator baseline across the full supervision budget range, down to $10\%$ root labels, and the supervision itself can be reallocated across the tree with little effect on accuracy.

\subsection{Paper Outline}

Section~\ref{sec:mergeable-sketches} opens with the prior literature and algebraic target: data-system mergeable summaries, ordered list homomorphisms, task-relative sufficient statistics, and the collapse theorem that relates those settings to C-TreePO. Section~\ref{sec:framework} formalizes the three local laws that operationalize that target for a learned summarizer, and Section~\ref{sec:theorems} states the four headline results: root preservation, preference-objective equivalence, the quantitative gap bound, and the audited certificate, with the full result ladder, the realizability layer, and the assumption-level stress tests in Appendix~\ref{app:theory-details}. Section~\ref{sec:estimation} turns the certificate into a deployable estimator: a design-based node-sampling audit with logged propensities, the object a survey methodologist would recognize. Section~\ref{sec:manifesto-deepdive} is the empirical centerpiece: the measurement setting, the document-label results (headline parity with the frontier ensemble, chunk invariance, the alternating prompt ladder and its joint-training failure mode), the quasi-sentence gold-label results (two label systems, trained substrates on identical bundles, gold-seeded summarization), and the lessons the corpus teaches the framework. Section~\ref{sec:discussion} states the gap-bound decomposition in plain language and locates the neighboring methods. The controlled validations and supporting material live in the appendices: the Markov changepoint interlude, neural-operator primer, and walk-through (Appendices~\ref{sec:markov-interlude}--\ref{sec:markov-walkthrough}); HyperLogLog parity and the classical-sketch suite (Appendices~\ref{sec:hll-parity} and~\ref{app:classical-parity}); full proofs and the Lean crosswalk (Appendices~\ref{app:proofs} and~\ref{app:proof-artifacts}); the Benoit replication record and quasi-sentence grids (Appendices~\ref{app:benoit-replication}--\ref{app:gold-seeding}); and the \texttt{treepo} software (Appendix~\ref{app:treepo-software}).
